\scanpage{28}
\mspage{9}
avec $E_J\cap T$ cônique dans $T$, et la structure cônique
\deleted{induit} étant la structure produit par $I^J$ d'
une structure cônique sur $(\overline T_J,Y_J)$, ayant
$\dot T_J$ comme partie-bord du cône sur $Y_J$\ldots

\deleted{\unsure{6}}

\textcircled{6.} \unsure{Relations|Relatifs} aux stratifications de type $I$
($I$ ens ordonné fini)
Une telle stratification (d'un espace $X$) est
la donnée d'une famille croissante de
sous-espaces fermés $X_i$. On dit que
la stratification est \uline{pré-cylindrique} si

\deleted{$\forall\unsure{i}$}
a) $\forall j'\unsure{\leftrightarrow|<}j$, $j$ successeur de $j'$,
$X_{j'}\hookrightarrow X_j$ est
\deleted{\unsure{une part.}} \inserted{bordant}

\inserted{b) \deleted{\unsure{Pour}}\inserted{Pour} $j$ fixé,
$j'$ variable parmi les prédécesseurs, les
$X_{j'}$ sont des parties \unsure{bordantes} transversaux de $X_j$}

\margin{i.e. munis des voisinages cylindriques transversaux.}

\deleted{b)}\inserted{c)}
$\forall i\unsure{\leftrightarrow|\le}j$, la structure de l'ens \inserted{$I_{i,j}$}
des $k$ tels que $i\le k\le j$ est celle de l'ens
\inserted{parties d'un ens} \deleted{simplexe} --- ce qui signifie que
si $S\inserted{(=S_{ij})}$ est l'ens. des prédécesseurs de $j$ dans
$I_{i,j}$ \inserted{i.e. des éléments maximaux de $I_{i,j}\setminus\{j\}$}
et si pour $k\in I_{i,j}$, on désigne
\deleted{\unsure{alors}} par $s_k$ \deleted{l'ens des} parties $s_k\subset S$ formée
de $j\in S_{\unsure{k}}$ tels que
$\unsure{j\ge s_k|j\ge k}$ (donc $s_k=s_{kj}$)
l'application $k\longmapsto s_k$ (qui est décroissante)
est un antiisomorphisme de $I_{ij}$ avec $\mathcal P(S)$
ordonné par inclusion. \deleted{$\unsure{I}$}

\margin{\uline{NB} Compte tenu de c), il suffit d'exiger dans
b) pour $j$ maximal dans $I$.}

\margin{\deleted{\unsure{dans un même cône}
\deleted{\unsure{des parties}}}}

\margin{On veut de plus \deleted{\unsure{que}}
que les \unsure{Inf} dans $I_{i,j}$ correspondent à
l'intersection dans $X_j$\ldots cf. (0) ci-dessous.}

Combinant b) et c), on voit que le
\margin{27}
